\documentclass{beamer}

\usepackage{graphicx}
\usepackage{calc}
\usefonttheme{serif}

\title{MiKeDa Bay Electrical}
\subtitle{JMCT}
\date{}

\usetheme{jmct}

%%%%%%%%%%%%%%%%%%%%%%%%%%%%%%%%%%%%%%%%
%%% Helpers (carried over from the nats deck)
%%%%%%%%%%%%%%%%%%%%%%%%%%%%%%%%%%%%%%%%

\newcommand{\textover}[3][l]{%
 % #1 is the alignment, default l
 % #2 is the text to be printed
 % #3 is the text for setting the width
 \makebox[\widthof{#3}][#1]{#2}%
 }

\newcommand{\blueit}[1]{%
  {\color{dark-lucid-blue}#1}%
}
\newcommand{\blueite}[1]{%
  \blueit{\emph{#1}}%
}

\newcommand{\myquote}[3]{
  ``#1''
  \vspace{3pt}
  \hrule
  \begin{flushright}
  --- \blueit{\emph{#2}}, \emph{#3}
  \end{flushright}
}

%%%%%%%%%%%%%%%%%%%%%%%%%%%%%%%%%%%%%%%%
%%% Helpers for this deck
%%%%%%%%%%%%%%%%%%%%%%%%%%%%%%%%%%%%%%%%

% A quantity with its unit: \val{48}{V}
\newcommand{\val}[2]{#1\,#2}

% A named piece of kit, styled consistently everywhere it appears
\newcommand{\kitname}[1]{{\color{dark-gray}#1}}

% Scaffolding marker. Every one of these is a hole the wiki needs to fill.
% `grep -n todo electrical.tex` to find what's left.
\newcommand{\todo}[1]{{\color{galois-blue}\textbf{[TODO]}}~\emph{#1}}

% A two-column spec block:
%   \begin{kit}
%     Inverter & 3000\,W \\
%   \end{kit}
\newenvironment{kit}%
  {\begin{center}\begin{tabular}{@{}l@{\hspace{1.2cm}}l@{}}}%
  {\end{tabular}\end{center}}

% The running outline. \signpost{n} highlights item n; \signpost{0} highlights
% nothing. Keeping it in one place means the section list only changes here.
\newcommand{\outitem}[3]{%
  \item \ifnum#1=#2 \blueit{#3}\else #3\fi
}
\newcommand{\signpost}[1]{%
  \begin{enumerate}
    \outitem{#1}{1}{Understand the current system}
    \outitem{#1}{2}{Talk Battery tradeoffs}
    \outitem{#1}{3}{Present 3 paths forward}
  \end{enumerate}
}

\begin{document}
	\frame {
		\titlepage
	}

%%%%%%%%%%%%%%%%%%%%%%%%%%%%%%%%%%%%%%%%
%%% Intro
%%%%%%%%%%%%%%%%%%%%%%%%%%%%%%%%%%%%%%%%



  \frame {
    \frametitle{Main takeaway}
    Batteries are likely past their healthy life. \blueite{something} will need to change.
  }
  \frame{
    \frametitle{Main takeaway}
    \begin{itemize}
      \item<2 -> Current batteries are not available to consumers. Papa Davis had a connection to the telecom industry that allowed him to get \blueite{ very special batteries}
      \item<3 -> These batteries traded longevity for \emph{ability to cycle}.
      \item<4 -> There are only two choices of battery type available to consumers like us.
    \end{itemize}
	}

%%%%%%%%%%%%%%%%%%%%%%%%%%%%%%%%%%%%%%%%
%%% Outline
%%%%%%%%%%%%%%%%%%%%%%%%%%%%%%%%%%%%%%%%

  %%%%%%%
  % Goals
  %

  %%%%%%%
  % Actual outline
  %
  \frame{
    \frametitle{Shape of things to come}
    \signpost{0}
  }

%%%%%%%%%%%%%%%%%%%%%%%%%%%%%%%%%%%%%%%%
%%% 1. The site and the loads
%%%%%%%%%%%%%%%%%%%%%%%%%%%%%%%%%%%%%%%%


  \frame{
    \frametitle{What we're working with}
    \begin{itemize}
      \item<2 -> 3 sets of solar panels
      \item<3 -> A bank of 12 \blueite{lead calcium} batteries
	  \item<4 -> the controllers to charge the batteries and the inverter to use the batteries
    \end{itemize}
  }
  \frame{
    \frametitle{The capacity}
    \blueite{As rate} the system has a capacity of 24kWh
	\begin{itemize}
      \item<2 -> This is a lie
      \item<3 -> This is a lie for two reasons:
    \end{itemize}
  }

  \frame{
    \frametitle{The capacity}
	\begin{itemize}
      \item<1 -> Only true when the batteries are within the `optimal' age
      \item<2 -> \blueite{Lead Calcium batteries do not like to be fully discharged!}
	  \item<3 -> Papa Davis took very good care not to discharge the batteries more than 30-50\%
	  \item<4 -> Combine that and age of the batteries, actual capacity likely closer to 10kWh.
    \end{itemize}

  }
  \frame{
    \frametitle{The Panels:}
    \begin{enumerate}
      \item<2 -> \blueite{Energy} --- how much you use in a day (sizes the bank
                 and the array)
      \item<3 -> \blueite{Power} --- how much you use at once (sizes the
                 inverter and the wire)
      \item<4 -> Conflating them is how you end up with a system that browns
                 out on a full battery
    \end{enumerate}
  }

%%%%%%%%%%%%%%%%%%%%%%%%%%%%%%%%%%%%%%%%
%%% 2. Generation
%%%%%%%%%%%%%%%%%%%%%%%%%%%%%%%%%%%%%%%%

  \frame{
    \frametitle{Sign post}
    \signpost{2}
  }

  \frame{
    \frametitle{Batteries: an intuition}
    \begin{itemize}
      \item<2 -> Lead Acid: can provide significant energy very quickly, but doesn't like to be fully drained: Think `car battery'.
      \item<3 -> Lithium: can take energy more quickly, and can be (more) fully discharged, but doesn't like to be used in the very-cold: Think mobile phone.
    \end{itemize}
  }
  \frame{
    \frametitle{Worth Repeating}
  	There is no `replace it with the same'. We are not able to get Lead Calcium batteries, their benefits/downsides are irrelevant to us.
  }
  \frame{
    \frametitle{General Tradeoff: Lead Acid}
	Lead Acid are \blueite{slightly} cheaper, and don't mind freezing. They don't last as long and give less \blueite{usable} capacity.
  }
  \frame{
    \frametitle{General Tradeoff: Lithium}
	Lead Acid are \blueite{slightly} more expensive, and can't be charged when freezing. They last years longer and give more capacity.
  }
  \frame{
    \frametitle{Cost Tradeoff}
	When accounting for the lifespan of the batteries, Lithium wins in cost-per-year.
  }
  

%%%%%%%%%%%%%%%%%%%%%%%%%%%%%%%%%%%%%%%%
%%% 3. Storage
%%%%%%%%%%%%%%%%%%%%%%%%%%%%%%%%%%%%%%%%

  \frame{
    \frametitle{Sign post}
    \signpost{3}
  }

  \frame{
    \frametitle{Possible Paths}
    \begin{itemize}
      \item<2 -> Replace Lead Calcium with Lead Acid, capacity stays basically the same, and we revisit in 8-ish years.
      \item<3 -> Replace Lead Calcium with Lithium, capacity increases and we have to bring the batteries home for winter. Will last 10-15 years/
	  \item<4 -> Split the two camps: one solar array powers our camp with new lithium system (because they take charge more quickly, this can actually work, you don't need as many solar panels for lithium).
    \end{itemize}
  }

\frame{
    \frametitle{Possible Paths}
    \begin{itemize}
      \item<2 -> The cost difference betweeen the cheapeast and most expensive is about \$5,000 CAD
      \item<3 -> Real question is whether we use this as reason to `split' the camps. 
	  \item<4 -> Benefit of split is that the current electrical setup probably has a few years left for the amount that grandpa's camp is used. So the system can be `unchanged' for grandpa's camp.
	  \item<5 -> Downside of keeping systems joint is that any future split would require purchasing more batteries anyway...
    \end{itemize}
  }


%%%%%%%%%%%%%%%%%%%%%%%%%%%%%%%%%%%%%%%%
%%% 4. Conversion
%%%%%%%%%%%%%%%%%%%%%%%%%%%%%%%%%%%%%%%%


\end{document}
